\section{Simulated customers and causal assistants need different experiments}
\label{exec20:pa-scope}

A bank considering a new servicing message may ask a language model two
different questions. How would a customer with a given history respond to the
message? Which features of the history justify a causal explanation of the
response? The first question invites a simulated respondent; the second asks
for assistance with a causal model. A convincing answer to either question
does not establish the other. This distinction determines what a public-data
prototype should test before an expensive bank experiment is justified.

Simmons and Hare review language models as subpopulation representative
models: systems intended to reproduce behavior or opinions of a specified
population. They organize prompting, fine-tuning, grounding data and deployment,
and distinguish direct comparison with observed responses, human assessment
and qualitative aggregate evaluation. Their five guiding criteria are fidelity,
necessity, robustness, sensitivity and fairness. In this use, a demographic
description is a proposed conditioning variable; it is not evidence that the
generated respondent represents a customer selected by the bank's actual
contact and response processes.\citep{exec20_population}

The selected May 2026 revision of Cai, Liu and Song instead manipulates inputs
to a causal assistant. It compares variable names, numerical information and
their removal, then studies instruction fine-tuning on pairwise discovery
tasks. These experiments concern the assistant's behavior under controlled
input conditions. They can guide a prototype evaluation without turning an
assistant's answer into customer-level intervention evidence.
\citep{exec20_attribution}

\subsection{A precise simulation can retain a population error}
\label{exec20:pa-population}
For a fixed customer segment, response definition and historical period, let
$p$ be the true probability of a specified binary response. Let $q$ be the
probability that a fixed language-model configuration generates that response.
Both probabilities refer to the same question, but to different populations
of outcomes. Suppose, for this project derivation, that $M$ independent
decoder draws have common success probability $q$, and let $\widehat q_M$
be their mean. Then
\begin{equation}
 E\{(\widehat q_M-p)^2\}
   =(q-p)^2+\frac{q(1-q)}{M}.
 \label{exec20:pa-mse}
\end{equation}
This follows by writing $\widehat q_M-p=(\widehat q_M-q)+(q-p)$;
the cross term has expectation zero. More draws reduce Monte Carlo error
around the model's own response law. They do not remove the discrepancy
between that law and the target population. If $p$ must itself be estimated
from a survey, uncertainty and selection error in that estimate are additional
obligations.\citep{exec20_population_attribution_reference}

For an exact assumed example, take $p=1/2$, $q=3/5$ and $M=10{,}000$.
The squared discrepancy is $1/100$, while the simulation variance is only
$3/125{,}000$. If a specified response earns USD 100 and each contact costs
USD 55, the model suggests USD 5 per contact while the stipulated population
law gives USD $-5$. This is a generated planning example, not an estimated
bank return. Moreover, a response forecast is not automatically an incremental
effect of contacting a customer: that decision also needs the relevant
no-contact potential outcome and identification assumptions.
\citep{exec20_population_attribution_reference}

The survey's distinction between consistency, pattern correspondence and
diversity is useful here. Producing a syntactically valid answer or a plausible
customer narrative is a weaker test than matching joint response patterns and
within-segment variation. Its distinction between apparent and actual model
sensitivity also cautions against treating a demographic prompt as a complete
description of the information on which the model relies.
\citep{exec20_population}

A proposal-stage comparison can therefore use a frozen public response
dataset, with its population, dates, response construct and reuse terms
documented. The comparison should include simple empirical or statistical
response models, preserve subgroup denominators, and evaluate held-out
responses and their joint variation. That establishes a bounded descriptive
task. A second, separately declared experiment must test an intervention
question under a known generated mechanism or a defensible public causal
design. Simulated respondents are useful for testing interfaces, eliciting
possible failure cases and comparing designs; their apparent realism supplies
no additional observed bank customers.

\subsection{Input ablations identify a controlled assistant experiment}
\label{exec20:pa-attribution}
In the attribution paper, context and model conditions are held fixed while
names associated with domain knowledge and numerical data are varied. Write
$a_{kd}$ for the correctness probability under a fully specified input
construction, where $k=1$ retains meaningful variable names, $d=1$ retains
numerical input, and zero denotes the corresponding removal. The probability
law must state whether it averages over tasks, decoder randomness, or both.
The paper's conditional and marginal attributions can then be written as
\begin{align}
 \mathrm{CAK}&=a_{11}-a_{01}, &
 \mathrm{CAD}&=a_{11}-a_{10}, \nonumber\\
 \mathrm{MAK}&=a_{10}-a_{00}, &
 \mathrm{MAD}&=a_{01}-a_{00}.
 \label{exec20:pa-factorial}
\end{align}
Here CAK and CAD retain the other input, while MAK and MAD compare against
the no-name, no-data baseline. The selected paper implements removal of
name-associated cues through arbitrary replacement names and removal of
data through prompts containing names alone. Its prompting discussion also
distinguishes wording that asks an assistant to ``suggest'' a relationship
from wording that may induce noncommittal responses.
\citep{exec20_attribution}

These are concrete changes to an interface. Replacing names does not delete
the model's learned parameters, remove every semantic cue in the context,
or prove that all other aspects of difficulty are unchanged. The defensible
experimental target is the effect of the declared prompt transformation on
assistant performance. Calling that transformation knowledge removal requires
additional construct-validity evidence. A refusal must remain an observed
outcome; a more encouraging prompt should not improve a score merely by
removing refusals from its denominator.

The equality $\mathrm{MAD}-\mathrm{MAK}=\mathrm{CAD}-\mathrm{CAK}$
follows algebraically from the four cells. It does not independently test
whether the prompt transformations isolate knowledge and data. To see what
the conditional comparisons count, define the project-derived interaction
\begin{equation}
 I=a_{11}-a_{10}-a_{01}+a_{00},\qquad
 \mathrm{CAK}+\mathrm{CAD}=(a_{11}-a_{00})+I.
 \label{exec20:pa-interaction}
\end{equation}
For assumed cell probabilities $(a_{11},a_{10},a_{01},a_{00})=
(9/10,3/5,2/5,3/10)$, the total gain is $3/5$, the interaction is $1/5$,
and the two conditional attributions sum to $4/5$. Adding them as independent
portions of improvement counts the interaction twice. The ordered paths
$\mathrm{CAK}+\mathrm{MAD}$ and $\mathrm{CAD}+\mathrm{MAK}$ each
equal the total gain. These exact values illustrate the estimand; they are
not measurements from the paper.\citep{exec20_population_attribution_reference}

For a banking assistant, a paired factorial design should freeze the model,
system instructions, permitted tools, decoding settings and response parser.
Each generating mechanism should appear under all declared input conditions,
with randomized ordering and separate conversation state. The test should
record correct, incorrect, ambiguous and refused outputs. Names should include
both informative and deliberately misleading descriptions, so that a familiar
story cannot substitute for the numerical assumptions that identify a graph.

\subsection{Residual independence must survive the instruction boundary}
\label{exec20:pa-residual}
The fine-tuning section supplies names, definitions and numerical results from
discovery algorithms in the instruction. It describes four pairwise scenarios
through residual correlation, including an undefined direction and the two
possible directed relations. It then teaches an assistant to use knowledge
when the numerical direction is undefined. This description raises a specific
interface question: does the actual numerical artifact test independence under
the required model, or only ordinary zero correlation? The inspected paper
does not settle the native implementation of that boundary.
\citep{exec20_attribution}

An exact project example shows why that distinction matters. Let independent
$X,U$ each take values $-1$ and $1$ with equal probability, and set $Y=2X+U$.
The forward population least-squares residual is $U$, independent of $X$.
The reverse slope and residual are
\begin{equation}
 b=\frac{\operatorname{Cov}(X,Y)}{\operatorname{Var}(Y)}=\frac25,
 \qquad R=X-bY=\frac{X-2U}{5}.
 \label{exec20:pa-reverse}
\end{equation}
Both $\operatorname{Cov}(X,U)$ and $\operatorname{Cov}(Y,R)$ are zero.
The latter is the ordinary least-squares orthogonality condition, not an
independence result. Direct enumeration gives
\begin{equation}
 E(R^2Y^2)=\frac9{25},\quad E(R^2)E(Y^2)=1,
 \quad \operatorname{Cov}(R^2,Y^2)=-\frac{16}{25}.
 \label{exec20:pa-dependence}
\end{equation}
Thus the reverse residual is dependent on its regressor despite zero linear
correlation. A correlation-only instruction can lose the relevant distinction.
This is not a claim that LiNGAM or the authors' executable artifact implements
that faulty rule; neither executable was run for this comparison.
\citep{exec20_population_attribution_reference}

The proposed adapter should consequently carry the numerical method, its
assumptions, the actual dependence test, sample size, uncertainty and refusal
state into the assistant's input. If the engine reports an unidentified
direction, a language-model suggestion can become a separately labeled
hypothesis or request for additional evidence. It must not silently turn that
state into an identified edge. A benchmark that rewards confident completion
of an undefined scenario needs a different acceptance interpretation from a
banking assistant that preserves uncertainty.

\subsection{Reported performance and a public qualification design}
\label{exec20:pa-design}
The attribution study uses nine datasets with fifteen column permutations
each, giving 135 task instances. Its TDR divides correctly predicted pairs by
true causal pairs; it also reports false discovery rate, structural Hamming
distance and F1. Those measures answer different questions, so high TDR alone
cannot certify that the assistant avoids additional false edges. The reverse
task describes keeping numerical values intact while changing the proposed
ordering or adjacency. Such a transformation is a useful input challenge,
but it does not itself generate observations from a new intervention law.
\citep{exec20_attribution}

For instruction fine-tuning, the paper reports 300 pairs producing 1,200
training instances and 106 pairs producing 424 testing instances. It names
Mistral-7B-v0.2 with a LoRA adapter in its training description. Table 4 reports
overall accuracy 0.90 for the fine-tuned model, compared with 0.82 for
GPT-4 turbo, while the undefined scenario scores 0.58 versus 0.66. The
comparison therefore supports a scoped reported gain, not superiority on
every scenario. Its later table uses ``Mixtral'' wording despite the main
training description; exact weight and code identity require resolution
before adoption. Equal scores on three additional tasks are not a universal
no-forgetting result. The paper itself leaves wider benchmark transfer,
finite-sample behavior, missing data and uncertainty-weighted evidence
combination as limitations or future work.\citep{exec20_attribution}

Repeated prompts quantify one kind of instability, not new independent
domains. Under an illustrative law in which all fifteen outcomes within each
dataset are identical and the nine datasets are independent, 135 calls contain
nine independent units. This extreme is an authored example, not an estimate
of the paper's dependence. A proposed comparison should preserve dataset and
mechanism groups, hold out causal pairs and generating mechanisms, and report
both within-task variability and variation across tasks.
\citep{exec20_population_attribution_reference}

The first prototype can be inexpensive. Its generated mechanisms should
include known directions, observational ambiguity, latent-confounding
violations, and the residual witness above. Its baselines should include a
numerical engine alone, a deterministic translation of that engine's result,
an untuned assistant and the proposed tuned assistant. Success requires
faithful preservation of numerical evidence and uncertainty, alongside any
improvement in explanation or discovery; agreement with a convenient answer
template alone is insufficient. Execution traces should separate the numerical
result, its transformation into an instruction and the assistant's final
claim, making a failed boundary localizable.

Public population-response evaluation and generated causal tasks establish
different merits. Together they can justify whether further prototype work is
worth funding, while keeping actual bank response truth, institutional review
and investment acceptance as later decisions. Neither a large synthetic
respondent pool nor a high assistant score replaces the banking question's
population, information clock, causal assumptions and uncertainty contract.
